\section{Persistent marking and recursive comparison}\label{sec:comparison}

The ordinary syntax is enlarged only for the purpose of defining comparison.
An ordinary head $f$ retains its declared base arguments
$\vec s=(s_1,\ldots,s_r)$. A marked head has the form
\[
 P_\tau=f^\star_\tau(\vec s;\vec e).
\]
Its base arguments still have their declared types. Its finite list of
extras $\vec e$ has arbitrary recorded types, and its output annotation
$\tau$ can differ from the original result type of $f$. A retyping $P_\rho$
changes only that output annotation. Variables and lambdas are retained,
and their subterms may contain marks. An \emph{ordinary} term has no marked
node anywhere. Every judgment below compares terms of equal outer type.

\begin{definition}[Persistent partial marking]
Define $\mk(x)$ to be undefined,
$\mk(\lambda x.s)=\lambda x.\mk(s)$ when the body mark exists, and
\[
 \mk(f(\vec s))=f^\star(\vec s;\varnothing),\qquad
 \mk(f^\star_\tau(\vec s;\vec e))=f^\star_\tau(\vec s;\vec e).
\]
Write $\headmark(s)$ for total top marking, which leaves a lambda spine
ending in a variable unchanged and otherwise agrees with $\mk$.
Thus $(\headmark)^2=\headmark$.
\end{definition}

\begin{definition}[Original-spine application]\label{def:application}
For $q:\rho_1\to\cdots\to\rho_n\to\tau$, the partial expression
$\ap q{a_1,\ldots,a_n}$ is defined as follows. An empty tuple returns $q$.
Consume its original initial lambdas, assigning the supplied arguments
simultaneously to their bound variables. If all arguments are consumed,
return the substituted residual term. If arguments remain, inspect the
original residual head. An ordinary or marked functional head is retained,
marked if necessary, its old arguments are substituted, and all remaining
arguments are appended to its extras at once. An original variable head
blocks the operation, even when its substituted value is functional.
\end{definition}

This last qualification is essential. With $\theta=o\to o$, a functional
head $h:\theta$, and $y:o$, the full call
$\ap{(\lambda x:\theta.x)}{h,y}$ is undefined, whereas the two successive
unary calls return $h^\star(;y)$. A successful full call does commute
with capture-avoiding substitution. This follows by following its original
lambda spine: the only branch that could inspect a substituted variable
is excluded by the success assumption.

\begin{definition}[Occurrence multiset extension]\label{def:multiset}
For a weak relation $W$ and a strict relation $S$, write
$A>_{\mathrm{mul}}^{W,S}B$ when the source occurrences split into
$A_0\uplus A_1$ with $A_1\ne\varnothing$, the target occurrences split
into $B_0\uplus B_1$, $A_0$ and $B_0$ are bijectively matched by
equal-typed $W$ judgments, and every occurrence of $B_1$ is $S$-below
some occurrence of $A_1$. A source occurrence may supply several new
occurrences; it may also supply none. Equal terms remain distinct
occurrences.
\end{definition}

Define $\succeq_\star,\succ_\star$ simultaneously as the least relations
closed under the following clauses. Every derivation is finite. In the
clauses concerning a marked owner, $P=f^\star(\vec s;\vec e)$; all copies
are retyped to their indicated types. Equivalent heads are denoted by
$f\approx g$, and strict head precedence by $f>g$.
\begin{enumerate}[label=\textsc{\arabic*.}]
\item[\textsc{Var}] $\W{x}{x}$.
\item[\textsc{Abs}] From $\W{s}{t}$ derive
$\W{\lambda x.s}{\lambda x.t}$, opening both binders with the same fresh name.
\item[\textsc{Fun}] $\W{f(\vec s)}{g(\vec t)}$ if the heads are equivalent,
their result types agree, and a permutation of $\vec t$ is pointwise
weakly matched with $\vec s$.
\item[\textsc{Put}] From $\W{f^\star(\vec s;\varnothing)}{t}$ derive
$\W{f(\vec s)}{t}$.
\item[\textsc{Select}] If $q$ is a base or extra argument of $P$, has type
$\rho_1\to\cdots\to\rho_n\to\tau$, and
\[
 \ap q{P_{\rho_1},\ldots,P_{\rho_n}}=r,\qquad \W r t,
\]
then $\W{P_\tau}{t}$. The tuple must succeed as one full call.
\item[\textsc{FAbs}] If
$\W{f^\star_\tau(\vec s;\vec e,x)}{t}$ after fresh opening at
$x:\sigma$, then
$\W{P_{\sigma\to\tau}}{\lambda x:\sigma.t}$.
\item[\textsc{Copy}] If $f>g$ and $\W{P_{\operatorname{type}(t_i)}}{t_i}$
for every target base argument, then $\W{P_{\operatorname{res}(g)}}{g(\vec t)}$.
\item[\textsc{Stat}] The same conclusion holds if $f\approx g$,
$\vec s>_{\mathrm{mul}}^{\succeq_\star,\succ_\star}\vec t$, and
all the same whole-source weak premises hold.
\end{enumerate}
The strict rule is
\[
 \frac{\mk(s)=s'\qquad \W{s'}{t}}{\St{s}{t}}.
 \tag{Strict}
\]
The base list, including its multiplicities, is used in \textsc{Stat}.
Extras are selectable but are absent from that multiset test.

\begin{remark}[Why persistence is material]
Let $Z<H<K<F$ be ground ranked heads of arities $0,0,1,2$, let $k=K(H)$,
and put
\[
 P=(\lambda x.F(x,Z))k,\qquad T=F(k,Z),\qquad U=F(H,H).
\]
Marking $P$ gives $q=@^\star(\lambda x.F(x,Z),k)$. Selection gives
$\W q H$, and persistence therefore gives $\St q H$. In
$F^\star(q,Z)$ remove both base occurrences and replace them by two
$H$ occurrences, each supplied strictly by $q$. The whole-source premises
to $H$ follow by selection of $q$. Thus \textsc{Stat} and \textsc{Put}
give $\W{F(q,Z)}U$. Selecting the lambda argument of $q$ with one
self-copy gives this very term, so $\St P U$.
This derivation uses an already-marked source in a strict multiset premise.
A definition that makes such a source unmarkable cannot use this derivation.
Our proofs below are for the persistent definition just given.
\end{remark}

\section{Structural laws, confinement and original reductions}
\label{sec:stability}

\begin{lemma}[Targets, inclusion and reflexivity]\label{lem:elementary}
Every target of $\succeq_\star$ or $\succ_\star$ is ordinary.
Strict comparison implies weak comparison. Every ordinary term is
weakly reflexive.
\end{lemma}
\begin{proof}
For the first assertion induct jointly over weak, strict and occurrence
list derivations. The variable target is ordinary. \textsc{Abs} and
\textsc{Fun} use their child targets; \textsc{Put} and \textsc{Select}
retain a premise target. \textsc{Copy} and \textsc{Stat} have an ordinary
target head and weak premises to every target argument. \textsc{FAbs}
binds its ordinary body target. The strict rule inherits its weak target.

If $\mk(s)=s'$ and $\W{s'}t$, peel the common initial lambdas.
At an ordinary head \textsc{Put} yields $\W s t$; at a marked head
$s'=s$. A variable cannot be marked. Lifting by \textsc{Abs} gives
strict inclusion. Finally structural induction on an ordinary term,
using \textsc{Var}, \textsc{Abs}, and identity matching in \textsc{Fun},
proves reflexivity.
\end{proof}

\begin{lemma}[Substitution and contexts]\label{lem:structural-stability}
The pair on ordinary endpoints is preserved by ordinary type-preserving
substitution and by every original well-typed context.
\end{lemma}
\begin{proof}
A successful mark commutes with substitution, including its marked case:
$\mk(s\theta)=\mk(s)\theta$. The equality follows by induction down
the initial lambda spine; substitution changes neither an ordinary head
nor a marked head. Successful full application also commutes with
substitution, by Definition~\ref{def:application}.

Induct simultaneously over weak, strict, weak-list matching and strict
replacement derivations. \textsc{Var} uses ordinary reflexivity of
$\theta x$. Under a binder use the lifted substitution. Head rules
transport their argument witnesses. \textsc{Select} uses the successful
application identity and the identity between substituted self-copies
and self-copies of the substituted owner. \textsc{FAbs} uses fresh
opening. \textsc{Strict} uses mark covariance. Applying a substitution
to a removed list preserves its nonempty set of occurrences even when
different terms acquire identical substituted values.

Weak context closure follows by induction on the context, using
\textsc{Abs} and \textsc{Fun}. For a strict hole under an ordinary head,
mark that head and use \textsc{Stat}: retain every unchanged occurrence
and replace the distinguished source argument strictly. Select the
corresponding source argument to supply each whole-source side premise,
using reflexivity or strict inclusion as appropriate. Under a lambda
lift the successful mark and its weak proof. Literal application is
one of these ordinary heads, so both of its argument positions are covered.
\end{proof}

\begin{lemma}[Free-variable containment]\label{lem:fv}
If $\W s t$ or $\St s t$, then $\FV(t)\subseteq\FV(s)$.
\end{lemma}
\begin{proof}
Use the same mutual derivation induction. Head matching distributes target
variables among source arguments. Copy and Stat have whole-source premises
to every target argument. A Select result uses only variables of the
selected argument and its owner copies, all already in the owner.
Opening adds one fresh variable, and the target abstraction binds it.
Marking changes no free variable. These observations give each rule case.
\end{proof}

Composition requires more than the preceding structural inductions.
Appendix~\ref{app:composition} gives a conservative auxiliary relation
on marked targets and proves its ordinary-final cut. Its proof preserves
the genuine failure of unrestricted associativity of full application.

\begin{theorem}[Ordinary composition]\label{thm:composition}
On ordinary endpoints, weak comparison is transitive, and
\[
 \succeq_\star\,;\,\succ_\star\ \subseteq\ \succ_\star,\qquad
 \succ_\star\,;\,\succeq_\star\ \subseteq\ \succ_\star.
\]
Consequently strict comparison is transitive.
\end{theorem}
\begin{proof}
Embed both weak proofs in the auxiliary relation of
Appendix~\ref{app:composition} and apply Theorem~\ref{thm:cut}.
For strict-then-weak use $\mk(s)$ as the initial source in the same cut.
For weak-then-strict, Lemma~\ref{lem:extended-marking} transports the
first weak proof to the successfully marked endpoints; markability of
the intermediate term reflects to the source. Cut with the weak witness
for the second strict judgment and apply \textsc{Strict}.
Strict inclusion then gives strict transitivity.
\end{proof}

\begin{definition}[Confined pair]\label{def:confined}
For ordinary terms of equal type define
\[
 s\cw t\iff \W s t\ \land\ \Ty(t)\subseteq\Ty(s),\qquad
 s\cs t\iff \St s t\ \land\ \Ty(t)\subseteq\Ty(s).
\]
\end{definition}

\begin{theorem}[Stable pair]\label{thm:stable}
The weak confined relation is reflexive and transitive, its strict part
is included in it and transitive, and both mixed composition laws hold.
Both parts are stable under original typed contexts and ordinary
capture-avoiding substitutions.
\end{theorem}
\begin{proof}
Support inclusions compose. The relational assertions therefore follow
from Lemma~\ref{lem:elementary} and Theorem~\ref{thm:composition}.
For substitution use Lemma~\ref{lem:fv} and
Corollary~\ref{cor:typesupport-stability}; every image type introduced
on the target is also introduced at a source occurrence. A fixed context
adds the same types to both sides. Combine these facts with
Lemma~\ref{lem:structural-stability}. Well-foundedness is proved separately
in Section~\ref{sec:upper}, without assuming this composition theorem.
\end{proof}

\begin{theorem}[Exact ranked restriction]\label{thm:rpo}
On first-order $\mathcal F$-terms, including patterns with ground-type
variables, $\cs$ is exactly multiset recursive path ordering. Its weak
part consists of that strict ordering together with recursive multiset
equality.
\end{theorem}
\begin{proof}
Every type support is $\{o\}$. An unmarked first-order source uses only
\textsc{Var}, \textsc{Fun}, or \textsc{Put}; a marked one uses only
\textsc{Select}, \textsc{Copy}, or \textsc{Stat}. Every selectable argument
has ground type, so its tuple is empty. No extras, lambdas, or marks
inside base arguments are generated. In particular, strict multiset
premises always have ordinary source arguments.

Induct on a weak derivation, interpreting an ordinary source comparison
as weak RPO and a root-marked source comparison as strict RPO.
\textsc{Var} is equality. In \textsc{Fun}, either every matched pair
is recursively equal or at least one is strict, giving a strict multiset
decrease. \textsc{Put} includes strictness into weakness. Selection is
the subterm clause, Copy the strict-precedence clause, and Stat the
multiset clause. The whole-source side premises are precisely the usual
domination requirements.

Conversely use \textsc{Var}/\textsc{Fun} for recursive multiset equality
and induct on an RPO derivation. The subterm clause is empty-tuple
Select. Strict precedence uses Copy, and the same-head clause uses Stat.
Each target argument is reached from its weakly matched or strictly
dominating source occurrence by selection, using strict inclusion.
Thus every side premise is available. The two inductions take place
entirely in the ranked fragment.
\end{proof}

\begin{lemma}[Ordinary template grafting]\label{lem:grafting}
If $A$ is an ordinary template and $\W p q$, then
$\W{A[x:=p]}{A[x:=q]}$, even when $p$ contains marks.
\end{lemma}
\begin{proof}
Induct on the ordinary template. At its distinguished variable use the
given comparison; other variables use Var. Abstraction uses fresh
opening, and every ordinary head uses Fun with identity argument
matching. The target image $q$ is ordinary by
Lemma~\ref{lem:elementary}, so the target remains ordinary.
\end{proof}

\begin{theorem}[Original reduction orientation]\label{thm:orient}
Every beta, $D0$, $DS$, $R0$, and $RS$ reduction, including its context
and substitution instances, strictly decreases under $\cs$.
\end{theorem}
\begin{proof}
For a beta redex $(\lambda x:\sigma.A)B:\tau$, put
$P_\rho=@^\star_\rho(\lambda x.A,B)$. Select its second base argument
to obtain $\W{P_\sigma}B$. Grafting gives
$\W{A[x:=P_\sigma]}{A[x:=B]}$. Selecting the first base argument
with its one self-copy produces exactly the former term. This gives
the required marked weak proof and hence strictness.

For $D^\tau uAB$, the outer second argument is $B$, so Select reaches
$B$. To reach $A$, select the outer first argument with one self-copy;
the resulting marked inner application still has $A$ as its second
base argument, selected with an empty tuple. These comparisons work
for every $u$, hence in particular for the two case rules. The $R0$
rule similarly selects the outer base argument.

For the successor recursor put $H=o\to\tau\to\tau$,
$U_1=R^\tau(Sn)$, $U_2=U_1A$, $L=U_2B$ and $L'=R^\tau nAB$.
Use the marked families
\[
\begin{aligned}
 P_\rho&=@^\star_\rho(U_2,B),\\
 C_{2,\rho}&=@^\star_\rho(U_1,A;P_\tau),\\
 C_{1,\rho}&=@^\star_\rho(R^\tau,Sn;C_{2,H},C_{2,\tau}),\\
 C_{0,\rho}&=(R^\tau)^\star_\rho(;C_{1,o},C_{1,H},C_{1,\tau}).
\end{aligned}
\]
Three Select prefixes carry a weak proof from $C_{0,\tau}$ back to
$P_\tau$, selecting $U_2,U_1,R^\tau$ with respectively one, two and
three simultaneous self-copies at their actual types.
The extra chain from $C_{0,H}$ reaches $A$. The chain from $C_{0,o}$
reaches $Sn$ and then $n$. Here $\St{Sn}n$ is itself outer Select.
Context closure of this strict comparison gives $\St L{L'}$, whose
weak witness is $\W{P_\tau}{L'}$. Prefixing the extra chain therefore
gives $\W{C_{0,\tau}}{L'}$.

Two Copy clauses, using the native head precedence $R^\tau>@$, now
construct first $An$ at type $\tau\to\tau$ and then $AnL'$ at type
$\tau$. Prefixing the three outer selections gives
$\W{P_\tau}{AnL'}$, as required. No native recursor has been treated
as a ranked ternary symbol.

Each contractum uses only node, binder and application-suffix types
already in its redex. Hence every displayed strict comparison is
confined. The stable-pair theorem gives all context and substitution
instances.
\end{proof}
